\MFQTransition{仓位与收益计算与应用分析题组}
\begin{enumerate}[leftmargin=*]
  \item \textbf{口述概念：}purge 与 embargo 分别阻断什么信息路径？
  \item \textbf{概念：}为什么目标仓位不能直接用于收益计算？
  \item \textbf{笔试推导：}逐期重建本章 $0.02$ 毛收益、2 单位换手和 $0.016$ 净收益。
  \item \textbf{数值编程：}把持有期改为 2，解释仓位和换手变化。
  \item \textbf{编程：}加入借券失败与停牌成交比例，扩展失败状态而非删除样本。
  \item \textbf{研究判断：}一项研究用完整样本拟合标准化参数，并按未成交的目标仓位计算收益。分别说明两种做法改变了什么信息或持仓假设，并重新写出正确的计算顺序。
  \item \textbf{综合项目：}选择一组可用数据，估计价差并制定仓位规则。比较简单基线，计算部分成交和成本后的结果，再分析关系变化时会发生什么。
\end{enumerate}

\subsection*{基础衔接：独立计算}
若上述交易的两腿总成本为 0.2，净损益是多少？下一期 beta 改变是否改变本期已实现损益？
